\section{Growing widths and coefficient modules}\label{sec:factors}

Keep the fixed geometry and separation of \cref{thm:fixed-width},
but at stage $i$ allow an arbitrary field $K_i$ of characteristic $p$
and arbitrary spaces $W_{v,i}$, of total dimension $B_i>0$.
The homogeneous coefficient maps may also vary. Choose an algebraic
closure $\Omega_i$ of $K_i$ and set
$\overline W_i=(\bigoplus_vW_{v,i})\otimes_{K_i}\Omega_i$.
Let $\mathcal A_i$ be the unital $\Omega_i$-algebra generated by
the vertex projectors and the embedded coefficient maps $T_e$.

Take a composition series of the $\mathcal A_i$-module
$\overline W_i$. Write its simple-factor dimensions as
$d_{i,1},\ldots,d_{i,m_i}$. For $k\ge1$ put
\begin{equation}\label{eq:factor-mass}
 \alpha_i(k)=\frac1{B_i}\sum_{j:d_{i,j}>k}d_{i,j}.
\end{equation}
These are dimensions over $\Omega_i$, not over $\F_p$.
All statements below hold for any composition series. The
Jordan--H\"older theorem also makes this multiset independent of the
chosen series; extension of a field to its algebraic closure does
not change a finite matrix's rank, since rank is detected by minors.

\begin{theorem}[A bound from simple-factor mass]\label{thm:factors}
\notformalized
For each fixed $k\ge1$ and $0<\varepsilon<1/2$, one index
$I(k,\varepsilon)$ works for all fields, widths, and homogeneous
coefficient tuples: for all $i\ge I(k,\varepsilon)$,
\[
 r_i\le(1-\alpha_i(k))(1/2+\varepsilon)+\alpha_i(k).
\]
\end{theorem}
\begin{informalproof}
Every invariant subspace $N$ is graded by types:
$N=\bigoplus_v E_vN$, because the vertex projectors $E_v$ are in
$\mathcal A_i$, are orthogonal, and sum to the identity.
The grading passes to every invariant subquotient, and each $T_e$
still maps its source piece to its target piece. Thus the induced
check on a simple factor fits the same template, with some
nonnegative vertex dimension vector of total $d_{i,j}$.

Tensor the composition series with $\Omega_i[Q_i]$ and apply
\eqref{eq:filtered-kernel}. For the factors of total dimension at
most $k$, there are only finitely many possible vertex dimension
vectors. Take the maximum of the corresponding thresholds from
\cref{thm:fixed-width}. This maximum is independent of the field,
coefficients, and number of factors. Every such factor contributes
at most $(1/2+\varepsilon)d_{i,j}q_i$ to nullity. Each larger factor
contributes at most $d_{i,j}q_i$. Sum, divide by $B_iq_i$, and use
$\sum_jd_{i,j}=B_i$. Field extension preserves the original rate.
\end{informalproof}

\begin{corollary}[Tightness and excess]\label{cor:factor-excess}
\notformalized
If
\[
 \lim_{k\to\infty}\limsup_i\alpha_i(k)=0,
\]
then $\limsup_i r_i\le1/2$. Conversely, if
$\liminf_i r_i\ge1/2+\delta$ with $0<\delta\le1/2$, then for
every fixed $k$,
\[
 \liminf_i\alpha_i(k)\ge2\delta.
\]
A limsup excess gives only the corresponding limsup conclusion.
\end{corollary}
\begin{informalproof}
For fixed $k$, take limsup in \cref{thm:factors} and then let
$\varepsilon\downarrow0$. This gives
$\limsup r_i\le1/2+(1/2)\limsup\alpha_i(k)$; now let $k\to\infty$.
For the converse, rearrange the finite inequality as
\[
 \alpha_i(k)\ge\frac{r_i-1/2-\varepsilon}{1/2-\varepsilon},
\]
take liminf, and then let $\varepsilon\downarrow0$.
The same rearrangement with limsup proves the last assertion.
No growing cutoff is substituted into a fixed-cutoff threshold.
\end{informalproof}

\begin{corollary}[A common internal flag]\label{cor:flag-half}
\notformalized
Suppose at each stage all $n=|V|$ vertex spaces are identified with
one space $Z_i$ of arbitrary positive dimension, and the internal
maps $C_e$ preserve a common full flag after algebraic closure.
Then $\limsup_i r_i\le1/2$, uniformly in fields and dimensions.
Pairwise commutativity of the identified internal maps suffices for
the flag hypothesis.
\end{corollary}
\begin{informalproof}
The direct sums over types of successive internal flag spaces are
invariant under every projector and embedded coefficient map.
Their successive quotients have total dimension $n$. Refining them
to composition series gives $\alpha_i(n)=0$; apply \cref{thm:factors}.

For commuting operators over an algebraically closed field, an
eigenspace of one operator is invariant under the others. Successively
restricting to eigenspaces gives a common eigenvector. Its invariant
line and induction on the quotient give a full common flag.
This proves simultaneous triangularizability, not diagonalizability.
The embedded maps between different types need not commute.
\end{informalproof}

\paragraph{Why factors are taken after algebraic closure.}
There is a strict difference from measuring simple factors over the
prime field. Take two types with identity maps in both directions
and one extra forward map given by multiplication by a generator
$\zeta$ of $\F_{p^m}/\F_p$. Use voltages $1,1,a$ over
$\Gamma=\langle a\rangle$ with quotients $C_{p^i}$ from
$\mathcal G=\Z_p\times\Z_p$ and
$U_i=p^i\Z_p\times p^i\Z_p$. Paired ports are disconnected
because the two incidence type classes are disjoint.
Over $\F_p$ the coefficient algebra is
$\operatorname{Mat}_2(\F_{p^m})$: the identity maps give the
cross matrix units and the extra map generates all field scalars
in a diagonal corner. Its natural module is simple of dimension $2m$.
Over an algebraic closure the separable field extension splits into
$m$ copies of that field, so the module has $m$ factors of dimension
$2$. Hence $\alpha_i(2)=0$ even when $m$ grows without bound.
This is a comparison of hypotheses; bipartiteness already bounds
every code on this particular relation by half rate.

The excess conclusion is a construction requirement, not a claim
that large simple factors produce high-rate codes. The explicit
example in \cref{sec:inflation} has maximal large-factor mass and
nevertheless obeys an elementary all-code half bound.
